\documentclass[../../main.tex]{subfiles}
\begin{document}

{\bf [解答]}

まず数字の穴埋めをし，その後(イ)から(ホ)の証明を行う．

\paragraph{数字の穴埋め}

題意の式
\begin{align}
  x_n + y_n \sqrt{3} = (2+\sqrt{3})^n
\end{align}
において$n=2,3$とすると
\begin{align}
  \begin{cases}
    x_2 + y_2 \sqrt{3} = (2+\sqrt{3})^2 = 7 + 4\sqrt{3} \\
    x_3 + y_3 \sqrt{3} = (2+\sqrt{3})^3 = 26 + 15\sqrt{3}
  \end{cases}
  \quad \therefore
  (x_2, y_2) = (7, 4), \quad (x_3, y_3) = (26, 15)
\end{align}
である．

次に漸化式について
\begin{align}
  x_n + y_n \sqrt{3}
  & = (2+\sqrt{3})^n \\
  & = (x_{n-1}+y_{n-1}\sqrt{3})(2+\sqrt{3}) \\
  & = (2x_{n-1}+3y_{n-1})x_n + (x_{n-1}+2y_{n-1})\sqrt{3}
\end{align}
だから，
\begin{align}
  x_{n} = 2x_{n-1}+3y_{n-1},
  \quad
  y_{n} = x_{n-1}+2y_{n-1}
\end{align}
である．

\paragraph{(イ)の証明}

与えられた条件
\begin{align}
  a + b\sqrt{3} = c + d\sqrt{3}
\end{align}
を変形すると
\begin{align}
  (a-c) = (d-b)\sqrt{3}
\end{align}
である．以下題意を背理法で示す．
もし $d-b \neq 0$ と仮定すると，
\begin{align}
  \dfrac{a-c}{d-b} = \sqrt{3}
\end{align}
だが，左辺は有理数，右辺は無理数となって矛盾する．
よって背理法より $d-b=0$ で, この時 $a-c=0$ が従う．
よって示された．

\paragraph{(ロ)の証明}

$x'=2x-3y, y'=2y-x$とおく．$(x,y)$が
\begin{align}
  f(x,y) = x^2-3y^2 = 1
\end{align}
を満たす時に$(x',y')$もこれを満たすことを示せば良い．
$f(x',y')$を計算すると
\begin{align}
  f(x',y')
  &= (x')^2 - 3(y')^2  \\
  &= (2x-3y)^2 - 3(2y-x) ^2 \\
  &= (4x^2-12xy+9y^2) - (3x^2-12xy+12y^2) \\
  &= x^2 -3y^2 = f(x,y) = 1
\end{align}
だから題意は示された．

\paragraph{(ハ)}

$x'=2x-3y, y'=2y-x$に対して与えられた不等式
\begin{align}
  x > x' > 0, \quad
  y > y' \ge 0
\end{align}
を示す．両者を代入して
\begin{align}
  &
  \begin{cases}
    x > 2x-3y > 0 \\
    y > 2y-x \ge 0
  \end{cases} \\
  \iff
  &
  \begin{cases}
    \dfrac{x}{3} < y < \dfrac{2x}{3} \\
    \dfrac{x}{2} \le y < x
  \end{cases} \\
  iff
  & \dfrac{x}{2} \le y < \dfrac{2x}{3}
\end{align}
を示せば良い．$x,y>0$より両辺二乗して$x^2-3y^2=1$を代入して
\begin{align}
  &\dfrac{x^2}{4} \le y^2 < \dfrac{4x^2}{9} \\
  iff
  &\dfrac{x^2}{4} \le \dfrac{x^2-1}{3} < \dfrac{4x^2}{9} \\
  \iff
  & 9x^2 \le 12(x^2-1) < 16x^2 \\
  \iff
  & 4 \le x^2, -3 < x^2
\end{align}
である．ここで，$x=1$の時は$x^2-3y^2=1$から$y=0となって$y>0$に反して矛盾する．
よって$x \ge 2$であり，この不等式は成立する．よって示された．

\paragraph{(ニ)}

任意の正の解$(x,y)$ から出発して(ロ)における$(x',y')$を$(x,y)$とする．
このような操作を順次行って，これを$X_1,X_2,\cdots$, $Y_1,Y_2,\cdots$とすると，(ハ)より
\begin{align}
  \begin{cases}
    X_1 > X_2 > \cdots  \\
    Y_1 > Y_2 > \cdots
  \end{cases}
\end{align}
となる．仮に$Y_m=0$なる$m$がないとすると任意の$n$に対して$Y_n>0$となり，$Y_1$が有限であることに反して矛盾である．よって$Y_m=0$なる$m$があり，このとき$X_m^2-3Y_m^2=1$より$X_m=1$である．よって題意は示された．

\paragraph{(ホ)}

(ニ)から，任意の解$(x,y)$から出発して
\begin{align}
  x'-y'\sqrt{3} = (2-\sqrt{3})(x-y\sqrt{3})
\end{align}
の操作を繰り返すことで$(x,y)=(x_0,y_0)=(1,0)$に辿り着くから，$n=0,1,\cdots$を用いて
\begin{align}
  & x_0 - y_0\sqrt{3} = (2-\sqrt{3})^n(x-y\sqrt{3}) \\
  \therefore
  & x-y\sqrt{3} = (2-\sqrt{3})^{-n} = (2+\sqrt{3})^{n}
\end{align}
と書ける．よって題意は示された．

\newpage

\end{document}
